El informe detalla el análisis experimental de cuatro algoritmos de ordenamiento unidimensional y dos algoritmos de multiplicación de matrices cuadradas. El objetivo fundamental es ir más allá de la notación asintótica teórica para evaluar el rendimiento práctico de estas implementaciones sobre hardware de un computador o dispositivo. Al someter los algoritmos a diferentes distribuciones y volúmenes de datos, se busca relacionar directamente su complejidad matemática con el tiempo de ejecución del algoritmo (programa), demostrando que factores como las constantes , la gestión de memoria y los peores casos teóricos tienen un impacto crítico en el poder de cómputo. Para poder entender esto a continuacion se presenta el entorno experimental con el que se trabajo.